\chapter{Environment Setup Reference}
\index{environment setup}\index{tooling reference}%
This appendix consolidates the environment requirements scattered through the book into one
reference you can revisit before any chapter. Chapter~0 walks through each step with narrative
and screenshots-in-words; this appendix is the checklist and command summary. Everything assumes
Windows with PowerShell, a sandbox subscription, and the book's naming and tagging conventions.

\begin{notebox}
Install once, verify always. Run \texttt{az account show --output table} at the start of every
session and before any destructive command. Most lab accidents come from the wrong subscription,
not from the wrong command.
\end{notebox}

\section{Azure Subscription and Microsoft Entra ID}
\index{Azure subscription}\index{Microsoft Entra ID}%
The labs assume the management hierarchy from Chapter~0: one Microsoft Entra tenant, one
subscription used only for this book, and one resource group per chapter named
\texttt{rg-de-azbook-chNN-dev} so cleanup is a single \texttt{az group delete}.

\begin{itemize}
    \item \textbf{Free account}: \$200 credit for 30 days, selected services free for 12 months,
    and always-free quotas \cite{microsoft2025azureFreeAccount}. The credit is not a safety
    net---a forgotten dedicated pool or premium cluster can exhaust it in days.
    \item \textbf{Identity separation}: never run daily labs as Global Administrator. Use a normal
    user account with \textbf{Contributor} on the lab resource groups, MFA enabled, and a separate
    break-glass admin account \cite{microsoft2025entra,microsoft2025rbac}.
    \item \textbf{Data-plane roles}: several labs need data-plane RBAC in addition to Contributor,
    notably \texttt{Storage Blob Data Contributor} on ADLS Gen2 accounts and Cosmos DB data-plane
    roles for keyless access. Owner on the resource is never a substitute.
    \item \textbf{Automation identity}: create resource-group-scoped service principals only when a
    lab requires non-interactive runs (Chapters 18, 24); store the secret in Key Vault, never in
    source control.
\end{itemize}

\section{Azure CLI}
\index{Azure CLI!setup}%
The primary command-line tool for every chapter \cite{microsoft2025azureCliInstall}.

\begin{lstlisting}[language=bash,caption={Installing and verifying Azure CLI on Windows.},label={lst:appenv-azcli}]
winget install --exact --id Microsoft.AzureCLI
az version
az login
az account list --output table
az account set --subscription "00000000-0000-0000-0000-000000000000"
az account show --output table
\end{lstlisting}

Useful extensions used by later chapters install on demand with
\texttt{az extension add --name <name>} (for example \texttt{databricks}, \texttt{datafactory},
\texttt{kusto}). Azure Cloud Shell is the no-install fallback for locked-down machines
\cite{microsoft2025cloudShell}, but copy commands and outputs back into your lab notes---its file
system and network path are not your machine's.

\section{Azure PowerShell}
\index{Azure PowerShell!setup}%
The \texttt{Az} module family is the first-class alternative when an organization standardizes on
PowerShell automation \cite{microsoft2025azurePowerShell}.

\begin{lstlisting}[language=bash,caption={Installing the Az module and selecting a context.},label={lst:appenv-azps}]
Install-Module -Name Az -Scope CurrentUser -Repository PSGallery -Force
Connect-AzAccount
Set-AzContext -Subscription "00000000-0000-0000-0000-000000000000"
Get-AzContext | Select-Object Name, Subscription, Tenant
\end{lstlisting}

\begin{pitfall}
\textbf{Split contexts.} \texttt{az account set} configures Azure CLI; \texttt{Set-AzContext}
configures Azure PowerShell. They are independent and can silently point at different
subscriptions in the same terminal. Check both before destructive operations.
\end{pitfall}

\section{Python SDKs and Local Tooling}
\index{Python SDK!setup}\index{development environment}%
Chapters with Python code assume Python 3.11+ with a per-book virtual environment. The recurring
packages are:

\begin{center}
\footnotesize
\begin{tabular}{@{}p{4.6cm}p{8.6cm}@{}}
\toprule
\textbf{Package} & \textbf{Used for} \\
\midrule
\texttt{azure-identity} & DefaultAzureCredential sign-in for all SDK clients \\
\texttt{azure-storage-file-datalake} & ADLS Gen2 file-system operations (Ch. 1, 7) \\
\texttt{azure-cosmos} & Cosmos DB NoSQL reads/writes (Ch. 3) \\
\texttt{azure-eventhub} & Event Hubs producers and consumers (Ch. 13) \\
\texttt{azure-ai-ml} & Azure ML workspaces, pipelines, endpoints (Ch. 17--19) \\
\texttt{redis} & Cache-aside client for Azure Cache for Redis (Ch. 4) \\
\texttt{pyodbc} / \texttt{sqlalchemy} & Azure SQL and Synapse connectivity (Ch. 2, 5) \\
\texttt{faker}, \texttt{pyarrow} & Synthetic data generation and Parquet writing (Ch. 5--6) \\
\texttt{databricks-cli} / \texttt{databricks-sdk} & Workspace automation (Ch. 9, 12, 15, 24) \\
\texttt{dbt-core} + \texttt{dbt-fabric}/\texttt{dbt-databricks} & Transformations as code (Ch. 22, 24) \\
\bottomrule
\end{tabular}
\end{center}

Keep credentials out of code: \texttt{DefaultAzureCredential} picks up the \texttt{az login}
session locally and the managed identity in Azure, so the same code runs in both places.

\section{Bicep and Terraform}
\index{Bicep}\index{Terraform!setup}%
Chapter~24 introduces infrastructure as code; install both tools early so later chapters can show
declarative alternatives to portal clicks.

\begin{lstlisting}[language=bash,caption={Installing Bicep and Terraform on Windows.},label={lst:appenv-iac}]
az bicep install && az bicep version
winget install --exact --id Hashicorp.Terraform
terraform version
\end{lstlisting}

Bicep is Azure-native and needs no state file; Terraform is the cross-cloud constant and stores
state in an Azure Storage account backend with lease locking \cite{microsoft2025bicep,hashicorp2025azurerm}.
Chapter~24's lab wires both into GitHub Actions with OIDC workload identity---no stored secrets.

\section{Cost Guardrails}
\index{cost guardrails}\index{budget!setup}%
Guardrails are part of the environment, not an afterthought. Before Chapter~1, create a monthly
subscription budget with alerts at 25\%, 50\%, 75\%, and 90\% of a small amount
\cite{microsoft2025costBudgets}, tag every resource group with
\texttt{project}, \texttt{environment}, \texttt{chapter}, and \texttt{owner}, and rehearse
deletion once:

\begin{lstlisting}[language=bash,caption={The budget-and-cleanup habit every chapter assumes.},label={lst:appenv-budget}]
az consumption budget create --budget-name budget-de-azbook-monthly `
  --amount 25 --time-grain Monthly --category Cost `
  --time-period start-date=2026-08-01 end-date=2027-08-01

az group delete --name rg-de-azbook-chNN-dev --yes --no-wait
\end{lstlisting}

\begin{tipbox}
End every lab session with two measurements: the resource list in the chapter's resource group
(\texttt{az resource list --resource-group <rg> --output table}) and the current forecast in Cost
Management. Pause or delete anything that bills while idle---dedicated pools, Databricks clusters,
VMs---before closing the terminal.
\end{tipbox}

\section{Per-Chapter Tool Requirements}
\index{tool requirements by chapter}%
The table below lists what each chapter expects to be installed or provisioned \emph{before}
Monday. ``CLI'' means Azure CLI; ``Python'' means the virtual environment above.

\begin{center}
\footnotesize
\begin{tabular}{@{}p{0.9cm}p{5.4cm}p{6.9cm}@{}}
\toprule
\textbf{Ch.} & \textbf{Primary services} & \textbf{Tooling prerequisites} \\
\midrule
0 & Entra ID, Cost Management & CLI, Azure PowerShell, budget, MFA \\
1 & ADLS Gen2 & CLI; Python: \texttt{azure-storage-file-datalake}; AzCopy optional \\
2 & Azure SQL, PostgreSQL Flexible Server & CLI; Python: \texttt{pyodbc}; sqlcmd optional \\
3 & Cosmos DB for NoSQL & CLI; Python: \texttt{azure-cosmos} \\
4 & Azure Cache for Redis, API Management & CLI; Python: \texttt{redis}, \texttt{pyodbc} \\
5 & Synapse dedicated SQL pool & CLI (\texttt{az synapse}); Python: \texttt{faker}, \texttt{pyarrow} \\
6 & Synapse dedicated SQL pool (existing) & Chapter 5 pool; T-SQL client \\
7 & Synapse serverless SQL pool & Chapter 5 workspace; ADLS data from Ch. 1 \\
8 & Microsoft Purview, Data Share & Purview account; scan credentials; Ch. 1 lake \\
9 & Azure Databricks & CLI \texttt{databricks} extension; workspace; ABFS access \\
10 & Azure Data Factory & CLI \texttt{datafactory} extension; Ch. 1 lake; Ch. 5 pool \\
11 & Azure DMS, Self-Hosted IR, Event Grid & Source Azure SQL DB; DMS; contract JSON schemas \\
12 & Databricks, ADF, Logic Apps, Event Grid & Ch. 9 workspace; Teams webhook for notifications \\
13 & Event Hubs, IoT Hub & Python: \texttt{azure-eventhub} (asyncio); Ch. 1 lake for Capture \\
14 & Azure Stream Analytics & Ch. 13 hub; target Azure SQL DB \\
15 & Databricks Structured Streaming & Ch. 9 workspace; Ch. 13 hub Kafka endpoint \\
16 & Power BI, Synapse & Power BI Desktop; Ch. 5 pool; gateway if on-prem \\
17 & Synapse Spark, AML, ONNX & \texttt{synapseml}, \texttt{onnxmltools}; Ch. 5 pool \\
18 & Azure ML, Databricks & Python: \texttt{azure-ai-ml}; AML workspace + compute cluster \\
19 & AML managed endpoints, monitoring & Ch. 18 registered model; Application Insights \\
20 & AI Document Intelligence, AI Search, Functions & Azure Functions Core Tools; Python: \texttt{azure-ai-*} \\
21 & Private Endpoints, Key Vault, Bastion & VNet + VM for tests; Key Vault; Event Grid \\
22 & Purview, Azure Monitor, dbt & dbt installed; Log Analytics workspace \\
23 & DDM, Always Encrypted, Defender, Policy & Ch. 2/5 databases; Key Vault for CMK \\
24 & Terraform, Bicep, dbt, GitHub Actions & Terraform + Bicep; GitHub repo; OIDC app registration \\
\bottomrule
\end{tabular}
\end{center}

Chapters are cumulative by design: the Chapter 5 warehouse, Chapter 9 Databricks workspace, and
Chapter 13 event hub are reused downstream, so tearing them down early forces rebuilds later.
Pause rather than delete shared infrastructure until the part of the book that uses it is done.
